\chapter{Non-Standard Analysis}
\label{ch:nsa}

Hierarchical classification, as used throughout this thesis, requires a learning
objective that respects a strict priority among per-level losses: a mistake on a
coarse label must never be offset by an improvement on a fine one. Expressing that
requirement inside a differentiable model turns out to be impossible over the real
numbers, and possible --- by an elementary formula --- as soon as the coefficients
of the objective are allowed to be infinitesimal. This chapter develops the
mathematical language in which that statement can be made precise.

The chapter is organised as follows. Section~\ref{sec:nsa-origins} recalls where
infinitesimals come from, why the nineteenth century eliminated them, and how
Robinson restored them to rigorous standing. Section~\ref{sec:nsa-motivation} shows,
with a short self-contained proof, that no real-valued loss can encode a
lexicographic priority: the obstruction is a fact about the real numbers, not a
matter of tuning weights. Section~\ref{sec:alpha} introduces the field used here, the
Euclidean numbers of Alpha Theory, with the two results needed to reason inside it
safely --- the transfer principle and the standard part.
Section~\ref{sec:lexico} then turns a lexicographic problem with any number of
objectives into a single non-standard scalar program.

Section~\ref{sec:nsa-bridge} is where that machinery is spent. It differentiates the
program and shows that neither the non-standard gradient nor its standard part can
train a real network: the first leaves the real parameter space, the second keeps
only the highest-priority level. What survives is a real update direction obtained by
projection, and that direction is the starting point of Chapter~\ref{ch:methods},
which builds the operators, states their guarantees, and works out the adaptation
each model family requires.

% =====================================================================
\section{Origins and development}
\label{sec:nsa-origins}
% =====================================================================

Non-standard analysis is not a recent invention answering a recent question. It is
the rigorous settlement of a dispute about the foundations of the calculus that
remained open for roughly three centuries. Since the construction used in this
thesis is an axiomatic shortcut through that history, it is worth recalling what the
shortcut avoids.

\subsection{Infinitesimals before rigour}

The differential calculus of Leibniz and Newton was formulated in terms of
\emph{infinitesimal} quantities: magnitudes greater than zero yet smaller than every
positive real number. A derivative was a ratio of infinitesimal increments, an
integral a sum of infinitesimal contributions. The technique was extraordinarily
productive, and the physics and analysis of the eighteenth century were built with
it.

Its logical status, however, was contested from the outset. An infinitesimal was
treated as non-zero when it appeared in a denominator and as zero when it was
discarded at the end of a computation, and no account was given of how the same
object could do both. The objection was pressed most famously by Berkeley, who
described such quantities as ``ghosts of departed quantities''. No answer was
available within the mathematics of the period, because the real number system
itself had not yet been characterised.

\subsection{The nineteenth-century settlement}

The nineteenth century resolved the difficulty by removing its subject. In the work
of Cauchy and, decisively, of Weierstrass, the infinitesimal was replaced by the
$\varepsilon$--$\delta$ formulation of limits: every statement previously made about
infinitely small quantities was rewritten as a statement about arbitrarily small but
strictly positive real quantities, quantified appropriately. The resulting analysis
was rigorous, and the price --- the loss of the infinitesimal as an object one could
name and compute with --- was regarded as no price at all.

The obstruction of Section~\ref{sec:nsa-motivation} is a consequence of that
settlement. It is not an accident of notation: once the number system is fixed to be
$\mathbb{R}$, certain orderings simply cannot be represented, as
Proposition~\ref{prop:debreu} shows.

\subsection{Robinson and the rehabilitation of the infinitesimal}

The situation changed in 1961, when Abraham Robinson showed that infinitesimals
could be given a fully rigorous foundation after all, using tools from model theory
that had not existed in Berkeley's time or in Weierstrass's. The construction was
developed in his 1966 monograph \emph{Non-standard
Analysis}~\cite{robinson1996nonstandard}, which gives the field its name.

Robinson's central device is what is now called the \emph{transfer principle}
(Section~\ref{sec:transfer}): a systematic licence stating that every elementary
property true of the real numbers remains true of their non-standard extension. It
is this principle that makes non-standard analysis usable rather than merely
consistent, since it allows the classical results of calculus to be reused verbatim
in the extended setting instead of being reproved.

The cost of Robinson's approach is the machinery it requires. Constructing the
hyperreal field demands ultrafilters and a background in model theory, which is a
substantial investment for a reader whose interest is computational rather than
foundational. This is the motivation for the alternative axiomatic presentations
discussed in Section~\ref{sec:alpha}.

\subsection{A family of non-Archimedean systems}

Robinson's hyperreals are not the only number system containing infinitesimals, and
the idea in fact predates him. Levi-Civita had already constructed a non-Archimedean
ordered field in 1892~\cite{levicivita1892infiniti}, obtained by adjoining to the
reals formal power series in an infinitesimal indeterminate; the field that bears his
name remains one of the standard examples, and the representation of its elements
anticipates the canonical form of Section~\ref{sec:euclidean-structure}. Much later,
Conway's theory of surreal numbers~\cite{conway2000numbers} produced a
proper class containing simultaneously the reals, the ordinals, and a vast supply of
infinite and infinitesimal magnitudes.

A separate line of work, pursued by Sergeyev, approaches numerical infinities and
infinitesimals through the \emph{grossone} methodology, with an explicit emphasis on
computation~\cite{sergeyev2017numerical}. It pursues goals close to those of the
present chapter by a different foundational route.

The system adopted here is \emph{Alpha
Theory}~\cite{benci2019measure,benci2022nonstandard}, chosen for two reasons that
matter for a thesis of this kind. Its presentation is elementary and self-contained,
requiring no model theory; and its numbers admit a finite-length binary encoding
suitable for a computer~\cite{benci2020algorithmic}, which keeps the construction
compatible with a numerical implementation even though, as
Remark~\ref{rem:algorithmic} records, this thesis never needs to exercise that
capability.

% =====================================================================
\section{Why non-standard analysis is needed here}
\label{sec:nsa-motivation}
% =====================================================================

Hierarchical classification (HC) is not a single learning task but a family of
learning tasks bound together by a priority relation: a mistake at the coarse level
of the taxonomy is qualitatively more damaging than a mistake at a fine level, because
the coarse decision constrains every decision below it. The natural formal counterpart
of ``qualitatively more damaging'' is the \emph{lexicographic} order, in which no
amount of improvement on a lower-priority objective can compensate a deterioration on
a higher-priority one.

\subsection{Lexicographic multi-objective optimisation}

The lexicographic order is formalised as follows.

\begin{definition}[Lexicographic multi-objective optimisation problem]
\label{def:lmop}
Let $f_1,\dots,f_n : \Omega \to \mathbb{R}$ be objective functions endowed with a
strict priority ordering induced by the natural order of their indices. The
\emph{lexicographic multi-objective optimisation problem} (LMOP) is the sequence of
programs
\[
  \min_{x \in \Omega} f_1(x),
  \qquad
  \min_{x} f_i(x) \ \ \text{s.t.}\ \ f_j(x) = \bar{f}_j \ \ \forall j = 1,\dots,i-1,\ x \in \Omega,
\]
where $\bar{f}_j$ is the optimal value attained by $f_j$ in the preceding
programs~\cite{fiaschi2024informed}.
\end{definition}

The definition describes a cascade rather than a single program: objective $i$ is
minimised only over the set on which objectives $1,\dots,i-1$ are already optimal.
Nothing obtainable at level $i$ can therefore be traded against anything at a level
above it.

\subsection{Why a real scalarisation cannot work}
\label{sec:lexico-real}

The standard machine-learning toolbox handles multiple objectives by
\emph{scalarisation}: the per-level losses are combined into a single real number
through a convex combination with weights $w_1,\dots,w_n \in \mathbb{R}_{>0}$,
possibly time-varying, as done for instance by branch convolutional neural
networks~\cite{zhu2017bcnn}. This device cannot, in principle, encode a lexicographic
priority --- and the obstruction holds for \emph{every} choice of positive real
weights, not just poorly-tuned ones. Denote by $\le_{\mathrm{lex}}$ the lexicographic
order on $\mathbb{R}^{n}$. The following classical result is due to
Debreu~\cite{debreu1954representation}; it is stated here for $n=2$, which suffices,
and shows that the obstruction is intrinsic.

\begin{proposition}[Non-representability of the lexicographic order in $\mathbb{R}$]
\label{prop:debreu}
There exists no function $u : \mathbb{R}^{2} \to \mathbb{R}$ such that
$(a,b) <_{\mathrm{lex}} (a',b') \iff u(a,b) < u(a',b')$.
\end{proposition}

\begin{proof}
Suppose such a $u$ exists. For every $a \in \mathbb{R}$ we have
$(a,0) <_{\mathrm{lex}} (a,1)$, hence $u(a,0) < u(a,1)$; by density of $\mathbb{Q}$ in
$\mathbb{R}$ choose $q(a) \in \mathbb{Q}$ with $u(a,0) < q(a) < u(a,1)$. Let $a < a'$.
Then $(a,1) <_{\mathrm{lex}} (a',0)$, so $u(a,1) < u(a',0)$ and therefore
\[
  q(a) < u(a,1) < u(a',0) < q(a').
\]
Thus $a \mapsto q(a)$ is strictly increasing, hence injective, and yields an injection
$\mathbb{R} \hookrightarrow \mathbb{Q}$. This contradicts $|\mathbb{R}| > |\mathbb{Q}|$.
\end{proof}

The obstruction is one of cardinality, not of smoothness or continuity: no
requirement was placed on $u$ beyond being a function into $\mathbb{R}$. Enlarging
the model class --- a deeper network, a different parameterisation, a schedule on the
weights --- therefore cannot circumvent it.

\begin{corollary}[Positive real weights always allow compensation]
\label{cor:compensation}
Let $w_1, w_2 \in \mathbb{R}_{>0}$ and $u(a,b) = w_1 a + w_2 b$. Then for every pair
with $\Delta_1 := a' - a > 0$ there exist values of $\Delta_2 := b' - b$, namely all
$\Delta_2 < -(w_1/w_2)\Delta_1$, for which $u(a',b') < u(a,b)$ even though
$(a,b) <_{\mathrm{lex}} (a',b')$.
\end{corollary}

\begin{proof}
Immediate: $u(a',b') - u(a,b) = w_1\Delta_1 + w_2\Delta_2 < 0$ precisely when
$\Delta_2 < -(w_1/w_2)\Delta_1$, and $w_2 > 0$ makes this set non-empty.
\end{proof}

This is the formal statement of a phenomenon observed empirically in hierarchical
classifiers trained with weighted losses: a sufficiently large gain on a fine-grained
objective can, and does, buy a loss on a coarse-grained one. Fiaschi and
Cococcioni~\cite{fiaschi2024informed} document exactly this behaviour for the
time-varying scalarisation of B-CNNs~\cite{zhu2017bcnn}, where the accuracy on the
highest-priority level degrades at the epochs in which the weights are shifted.
Setting $w_2=0$ removes the compensation but also removes the second objective from
the optimisation altogether, so it is not a solution.

\subsection{Why infinitesimal weights do work}
\label{sec:lexico-nonstd}

Proposition~\ref{prop:debreu} rules out every \emph{real} coefficient pair. It does
not rule out coefficients drawn from a larger field, and
Section~\ref{sec:nsa-origins} has already established that such fields exist. We
anticipate two pieces of notation from Section~\ref{sec:alpha}:
$\mathbb{E} \supseteq \mathbb{R}$ denotes the ordered field constructed there, and
$\eta \in \mathbb{E}$ a fixed positive \emph{infinitesimal} --- a number smaller than
every $1/n$, $n \in \mathbb{N}$ (Definition~\ref{def:magnitudes}). Only this defining
property of $\eta$ is used below.

\begin{proposition}[Lexicographic representability in $\mathbb{E}$]
\label{prop:e-representation}
Let $U : \mathbb{R}^{2} \to \mathbb{E}$ be defined by $U(a,b) = a + b\eta$. Then
\[
  (a,b) <_{\mathrm{lex}} (a',b') \iff U(a,b) < U(a',b').
\]
\end{proposition}

\begin{proof}
Write $\Delta_1 = a' - a$ and $\Delta_2 = b' - b$, so that
$U(a',b') - U(a,b) = \Delta_1 + \Delta_2\eta$.

\emph{Case $\Delta_1 > 0$.} The quantity $\Delta_2 \eta$ is infinitesimal (trivially
if $\Delta_2=0$; otherwise because $\eta$ is). Since $\Delta_1$ is a positive real, the
Archimedean property of $\mathbb{R}$ gives $m \in \mathbb{N}$ with $1/m < \Delta_1$;
as $|\Delta_2\eta|$ is smaller than $1/m$ in particular, $|\Delta_2\eta| < \Delta_1$
and $\Delta_1 + \Delta_2\eta > 0$.

\emph{Case $\Delta_1 < 0$.} Symmetric, giving $\Delta_1 + \Delta_2\eta < 0$.

\emph{Case $\Delta_1 = 0$.} Then $U(a',b') - U(a,b) = \Delta_2\eta$, positive iff
$\Delta_2 > 0$, since $\eta > 0$.

The three cases reproduce exactly the definition of $<_{\mathrm{lex}}$.
\end{proof}

Propositions~\ref{prop:debreu} and~\ref{prop:e-representation} together are the
precise sense in which infinitesimals are not a convenience but a necessity: the
lexicographic priority \emph{cannot} be encoded by any real-valued aggregate, and
\emph{can} be encoded, faithfully and by an elementary formula, as soon as the
aggregate is allowed to live in a field containing infinitesimals. Making this
precise --- what $\mathbb{E}$ and $\eta$ actually are, and which calculus results
survive the move from $\mathbb{R}$ to $\mathbb{E}$ --- is the subject of the next
section; Section~\ref{sec:lexico} then generalises
Proposition~\ref{prop:e-representation} from two to an arbitrary number $n$ of
objectives.

% =====================================================================
\section{Non-Archimedean fields and Alpha Theory}
\label{sec:alpha}
% =====================================================================

\begin{axiom}[Axiom of Archimedes]
\label{ax:archimedes}
Let $F$ be a totally ordered field. Then
$\forall\, x,y \in F,\ 0 < x < y,\ \exists\, n \in \mathbb{N} : y < nx$.
\end{axiom}

The axiom states that no element of $F$ is so large --- nor any positive element so
small --- that it cannot be reached by repeatedly adding another one to itself. It is
called \emph{Archimedean} after Otto Stolz, since it appears as Axiom~V of
Archimedes' \emph{On the Sphere and Cylinder}~\cite{benci2022nonstandard}. A totally
ordered field that does not satisfy it is called \emph{non-Archimedean}.

\begin{remark}[The real line is not an option]
\label{rem:R-archimedean}
$\mathbb{R}$ is Archimedean; more strongly, it is, up to order-isomorphism, the
unique Dedekind-complete totally ordered field, and every Dedekind-complete totally
ordered field is Archimedean. Infinite and
infinitesimal quantities therefore cannot be added to $\mathbb{R}$ without giving up
completeness: one must move to a strictly larger, necessarily incomplete, ordered
field --- e.g.\ the Levi-Civita field~\cite{levicivita1892infiniti}, the hyperreal
numbers of Robinson~\cite{robinson1996nonstandard}, or the surreal numbers of
Conway~\cite{conway2000numbers}.
\end{remark}

\subsection{The field of Euclidean numbers}

Alpha Theory introduces the field it needs by postulating it.

\begin{axiom}
\label{ax:field}
There exists an ordered field $\mathbb{E} \supseteq \mathbb{R}$, whose elements are
called \emph{$\alpha$-Euclidean numbers}.
\end{axiom}

Since $\mathbb{E}$ is a field, its elements obey the ordinary rules of algebra
(commutativity, associativity, distributivity, existence of inverses); non-standard
expressions can therefore be simplified exactly as real ones.

\begin{definition}[Infinite, finite, infinitesimal and appreciable numbers]
\label{def:magnitudes}
Let $\xi \in \mathbb{E}$. Then
\begin{align*}
  \xi \text{ is \emph{infinite}}      &\iff \forall\, n \in \mathbb{N},\ |\xi| > n, \\
  \xi \text{ is \emph{finite}}        &\iff \exists\, n \in \mathbb{N},\ |\xi| < n, \\
  \xi \text{ is \emph{infinitesimal}} &\iff \forall\, n \in \mathbb{N},\ |\xi| < \tfrac{1}{n}, \\
  \xi \text{ is \emph{appreciable}}   &\iff \exists\, n \in \mathbb{N},\ \tfrac{1}{n} < |\xi| < n.
\end{align*}
Write $\mathbb{E}_{\mathrm{fin}} \subset \mathbb{E}$ for the set of finite elements.
\end{definition}

These are not four disjoint classes. Infinite and finite are complementary, every
infinitesimal is finite, and the appreciable numbers are exactly the finite
non-infinitesimal ones. So $0$ is infinitesimal and every non-zero real is
appreciable --- but not conversely: $1+\eta$ is appreciable, hence finite, and is not
real. The finite numbers are therefore a strictly larger set than $\mathbb{R}$. What
does hold is that each of them sits infinitely close to exactly one real,
\begin{equation}
  \label{eq:finite-decomposition}
  \xi \in \mathbb{E}_{\mathrm{fin}}
  \quad\Longrightarrow\quad
  \exists!\, r \in \mathbb{R} \ \ \text{and} \ \ \exists\, \zeta \ \text{infinitesimal}
  \ : \ \xi = r + \zeta ,
\end{equation}
giving the strict inclusions
$\mathbb{R} \subsetneq \mathbb{E}_{\mathrm{fin}} \subsetneq \mathbb{E}$. The map
$\xi \mapsto r$ is the standard part of Section~\ref{sec:standard-part}: it is well
defined because $r$ is unique, and it does real work --- rather than acting as the
identity --- precisely because $\mathbb{E}_{\mathrm{fin}}$ is bigger than
$\mathbb{R}$.

Alpha Theory postulates one specific infinite number, $\alpha$, via the notion of
\emph{numerosity}, a refinement of cardinality sensitive to strict inclusion.

\begin{axiom}
\label{ax:num}
There exists a function $\mathrm{num}$, defined on the countable sets of a suitable
superstructure over $\mathbb{N}$, taking values in $\mathbb{E}$, such that
$\mathrm{num}(A) = |A|$ for $A$ finite; $\mathrm{num}(A) < \mathrm{num}(B)$ whenever
$A \subset B$; $\mathrm{num}$ is additive and multiplicative with respect to
$\cup$ and $\times$; and $\alpha := \mathrm{num}(\mathbb{N})$.
\end{axiom}

The second property is what distinguishes numerosity from cardinality: a proper
subset is strictly less numerous than its superset, even when both are infinite.
Setting $\eta := \alpha^{-1}$, the following hold in Alpha
Theory~\cite{benci2022nonstandard}:
\[
  0 < \eta = \alpha^{-1} < \alpha^{0} = 1 < \alpha < \alpha+1,
  \qquad
  \alpha \cdot (\alpha + 2) = \alpha^{2} + 2\alpha.
\]
By Definition~\ref{def:magnitudes}, $\alpha$ is infinite and $\eta$ is
infinitesimal --- the object anticipated in Proposition~\ref{prop:e-representation}.
The second identity is worth pausing on: $\alpha$ obeys the ordinary rules of algebra.
Nothing collapses, as it would in cardinal arithmetic, where
$|\mathbb{N}|+1 = |\mathbb{N}|$, or in ordinal arithmetic, where
$1+\omega = \omega \ne \omega+1$. These two symbols, $\alpha$ and $\eta$, are used
with this meaning throughout the thesis.

\begin{remark}[What the axioms leave open]
Axioms~\ref{ax:field}--\ref{ax:num} do not pin down every property one might
attribute to $\alpha$ --- whether it is even or odd, for instance, or the value of
$\sin(\alpha)$, is independent of them, and further axioms are needed to settle
questions of this kind~\cite{benci2019measure,benci2022nonstandard}. None of what
follows in this thesis depends on resolving them.
\end{remark}

\subsection{The structure of a Euclidean number}
\label{sec:euclidean-structure}

The axioms above assert that $\mathbb{E}$ exists and fix two of its elements, but
they do not yet say what a Euclidean number \emph{looks like}. The following two
statements supply that, and are the point at which $\mathbb{E}$ becomes an object one
can write down and, in principle, store.

\begin{definition}[Monosemium]
\label{def:monosemium}
Any value $\xi = r\alpha^{k} \in \mathbb{E}$, with $r,k \in \mathbb{R}$, is called a
\emph{monosemium}~\cite{fiaschi2024informed}.
\end{definition}

A monosemium is thus a single power of $\alpha$ scaled by a real coefficient. Its
exponent $k$ fixes its \emph{magnitude}: for $k>0$ the monosemium is infinite, for
$k=0$ it is the real number $r$, and for $k<0$ it is infinitesimal. The infinitesimal
$\eta = \alpha^{-1}$ used throughout this thesis is the monosemium with $r=1$ and
$k=-1$.

\begin{theorem}[Canonical form of a Euclidean number]
\label{thm:canonical-form}
Every $\xi \in \mathbb{E}$ can be written as a sum of monosemia,
\begin{equation}
  \label{eq:canonical-form}
  \xi = \sum_{i=1}^{\infty} r_i\,\alpha^{k_i},
  \qquad r_i \in \mathbb{R},\ k_i \in \mathbb{R},
  \qquad k_1 > k_2 > k_3 > \cdots,
\end{equation}
that is, as a series of monosemia with real coefficients and strictly decreasing real
exponents~\cite{fiaschi2024informed,benci2022nonstandard}.\footnote{The statement is
reproduced in the form used by the sources cited, with the monosemia indexed by
$\mathbb{N}$. In full generality $\mathbb{E}$ is a field of Hahn series over
$\mathbb{R}$ with real exponents: the support $\{k_i\}$ of $\xi$ is required only to
be a \emph{well-ordered} subset of $\mathbb{R}$, whose order type may exceed $\omega$,
and \eqref{eq:canonical-form} is then a formal Hahn sum rather than a limit in a
topology on $\mathbb{E}$. The distinction is immaterial here: every Euclidean number
this thesis writes down has finite support, and for finitely supported elements the
two formulations coincide.}
\end{theorem}

\begin{proof}
See~\cite{benci2019measure,benci2022nonstandard}.
\end{proof}

Two consequences of \eqref{eq:canonical-form} matter later. First, the ordering of
$\mathbb{E}$ is governed by the leading term: since the exponents decrease, the sign
of a non-zero $\xi$ is the sign of $r_1$, and comparisons between Euclidean numbers
reduce to lexicographic comparisons of their coefficient sequences. This is the
structural reason why an $\eta$-weighted sum encodes a lexicographic order, and it is
what Proposition~\ref{prop:e-representation} exploits in the two-objective case.
Second, a Euclidean number that involves only finitely many monosemia is completely
described by finitely many real pairs $(r_i,k_i)$, which is precisely the
representation exploited by the algorithmic numbers of
Remark~\ref{rem:algorithmic}.

\begin{remark}[Machine representation]
\label{rem:algorithmic}
Alpha-Theoretic numbers admit a finite-length binary
encoding --- \emph{algorithmic numbers}~\cite{benci2020algorithmic} --- with
constant-cost arithmetic, which is what makes Alpha Theory usable for actual
computation and not only for derivation. This thesis, however, never needs to
instantiate one: as shown in Chapter~\ref{ch:methods}, every non-standard quantity
introduced by the construction is eliminated analytically before any arithmetic is
performed on hardware.
\end{remark}

\subsection{The \texorpdfstring{$\alpha$}{alpha}-limit and the transfer principle}
\label{sec:transfer}

\begin{axiom}
\label{ax:alphalimit}
Every sequence $\varphi : \mathbb{N} \to \mathbb{R}$ has a unique
\emph{$\alpha$-limit}, denoted $\lim_{n \uparrow \alpha} \varphi(n)$, such that:
(1) every $\xi \in \mathbb{E}$ is the $\alpha$-limit of some real sequence;
(2) $\lim_{n \uparrow \alpha} n = \alpha$;
(3) if $\varphi(n) \ge \psi(n)$ eventually, then
$\lim_{n \uparrow \alpha} \varphi(n) \ge \lim_{n \uparrow \alpha} \psi(n)$;
(4) the $\alpha$-limit is linear and multiplicative:
$\lim_{n \uparrow \alpha}(\varphi(n)+\psi(n)) = \lim_{n \uparrow \alpha}\varphi(n) + \lim_{n \uparrow \alpha}\psi(n)$,
and likewise for products.
\end{axiom}

The notation $n \uparrow \alpha$ distinguishes the $\alpha$-limit from the ordinary
(Cauchy) limit $n \to \alpha$. The substantive content of
Axiom~\ref{ax:alphalimit} is property (1): \emph{every} sequence admits an
$\alpha$-limit, with no convergence hypothesis --- a model satisfying
Axioms~\ref{ax:field}--\ref{ax:alphalimit} can be built within
ZFC~\cite{benci2019measure,benci2022nonstandard}.

The $\alpha$-limit is what allows real functions to be carried into $\mathbb{E}$.

\begin{theorem}[Extension of a real function]
\label{thm:star-extension}
Every function $f:\mathbb{R}\to\mathbb{R}$ can be coherently and uniquely extended to
a function $f^{*}:\mathbb{E}\to\mathbb{E}$ such that $f^{*}(x) = f(x)$ for all
$x \in \mathbb{R}$. The extension is given by
\begin{equation}
  \label{eq:star-extension}
  f^{*}\!\left(\lim_{n \uparrow \alpha} \varphi(n)\right)
  = \lim_{n \uparrow \alpha} f\big(\varphi(n)\big),
\end{equation}
independently of the sequence $\varphi$ chosen to represent the
argument~\cite{fiaschi2024informed,benci2022nonstandard}.
\end{theorem}

\begin{proof}
See~\cite{benci2019measure,benci2022nonstandard}.
\end{proof}

The star is dropped when no ambiguity arises. With the extension in place, the
central licence of non-standard analysis can be stated --- but the word
\emph{elementary} has to be given a meaning first, since the principle is false
without the restriction it names.

Fix the superstructure $V(\mathbb{R}) = \bigcup_{m \ge 0} V_m$ over $\mathbb{R}$,
where $V_0 = \mathbb{R}$ and $V_{m+1} = V_m \cup \mathcal{P}(V_m)$, so that
$V(\mathbb{R})$ contains $\mathbb{R}$, every subset of $\mathbb{R}$, every relation
and function on $\mathbb{R}$, and so on upwards. Let $\mathcal{L}$ be the first-order
language with equality, the binary relation $\in$, and a constant symbol for each
element of $V(\mathbb{R})$.

\begin{definition}[Elementary formula]
\label{def:elementary}
A formula of $\mathcal{L}$ is \emph{elementary}, or \emph{bounded}, when every
quantifier occurring in it is restricted to a named set, i.e.\ appears in one of the
forms $\forall\, x \in A$ or $\exists\, x \in A$ with $A$ a constant symbol.
\end{definition}

\begin{theorem}[Transfer principle]
\label{thm:transfer}
Let $\sigma(x_1,\dots,x_k)$ be an elementary formula of $\mathcal{L}$ and let
$A_1,\dots,A_k \in V(\mathbb{R})$. Then
\begin{equation}
  \label{eq:transfer}
  \sigma(A_1,\dots,A_k) \iff \sigma({}^{*}A_1,\dots,{}^{*}A_k),
\end{equation}
where ${}^{*}A$ denotes the non-standard extension of $A$.
\end{theorem}

\begin{proof}
See~\cite{robinson1996nonstandard} for the model-theoretic formulation
and~\cite{benci2019measure,benci2022nonstandard} for the Alpha-Theoretic one.
\end{proof}

\begin{example}[Transferring a trigonometric identity]
To transfer $\sin(2x)=2\cos x\sin x$ to $\mathbb{E}$: let $\xi=\lim_{n\uparrow\alpha}
\varphi(n)$ for some real sequence $\varphi$. The identity holds pointwise for every
real $\varphi(n)$; taking the $\alpha$-limit of both sides and using linearity
(Axiom~\ref{ax:alphalimit}(4)) gives $\sin(2\xi)=2\cos\xi\sin\xi$. The same scheme
transfers any identity between real functions.
\end{example}

Boundedness is what makes the principle true. Dedekind completeness quantifies over
\emph{all} bounded subsets of $\mathbb{R}$, is therefore not elementary, and does not
carry over to $\mathbb{E}$ --- in agreement with Remark~\ref{rem:R-archimedean}. The
opposite case is the one that misleads.

\begin{remark}[What transfer does and does not say]
\label{rem:transfer-scope}
The Axiom of Archimedes (Axiom~\ref{ax:archimedes}) \emph{is} elementary, since its
quantifier over $n$ is bounded by $\mathbb{N}$, and it therefore transfers: for all
$x,y \in \mathbb{E}$ with $0 < x < y$ there exists $n \in {}^{*}\mathbb{N}$ with
$y < nx$. This is consistent with $\mathbb{E}$ being non-Archimedean, because the
witness is drawn from ${}^{*}\mathbb{N}$ and may be infinite. Taking $x = \eta$ and
$y = 1$, no $n \in \mathbb{N}$ satisfies $1 < n\eta$, whereas $n = \alpha + 1$ does,
since $(\alpha+1)\eta = 1 + \eta > 1$. Transfer preserves the \emph{form} of a
statement, not the standardness of the objects witnessing it; reading the transferred
conclusion back as a statement about $\mathbb{N}$ is what must be avoided.
\end{remark}

What this thesis takes from Theorem~\ref{thm:transfer} is the licence to reuse the
classical differentiation rules --- the ones backpropagation implements --- as valid
statements about $\mathbb{E}$. That is what makes the gradient of
Section~\ref{sec:nsa-bridge} legitimate, and it is all that is needed there: the
objective built in that section is a finite $\mathbb{E}$-linear combination of
standard real functions, so no function is ever star-extended. Full transfer is
stated because it is the defining feature of non-standard analysis and because it
fixes exactly what may and may not be carried across.

\subsection{The standard part}
\label{sec:standard-part}

\begin{theorem}[Standard part]
\label{thm:standard-part}
Let $\mathbb{E}_{\mathrm{fin}}$ be the set of finite elements of $\mathbb{E}$
(Definition~\ref{def:magnitudes}). There exists a function
$\mathrm{st} : \mathbb{E}_{\mathrm{fin}} \to \mathbb{R}$ such
that $x \sim y \implies \mathrm{st}(x) = \mathrm{st}(y)$ and
$\mathrm{st}(\mathrm{st}(x)) = \mathrm{st}(x)$, where $x \sim y$ means $x-y$ is
infinitesimal. It is the map $\xi \mapsto r$ of the decomposition
\eqref{eq:finite-decomposition}, and it is surjective onto $\mathbb{R}$ and constant
on each class of $\sim$.
\end{theorem}

\begin{proof}
See~\cite{benci2019measure}.
\end{proof}

Informally, $\mathrm{st}$ maps a non-infinite non-standard number to the unique real
number infinitely close to it --- e.g.\ $\mathrm{st}(2.1-5\eta)=2.1$ and
$\mathrm{st}(-3)=-3$. In terms of the canonical form
\eqref{eq:canonical-form}, $\mathrm{st}$ retains the monosemium of exponent $0$ and
discards every infinitesimal term. It is the canonical way of extracting a real
answer from a non-standard computation.

What $\mathrm{st}$ deletes matters as much as what it keeps.
Section~\ref{ssec:nsa-why-not} shows that applying it to the gradient of the
non-standard objective deletes every level below the first, which is why this thesis
recovers a real update direction by projection instead.

% =====================================================================
\section{From LMOP to a non-standard scalar program}
\label{sec:lexico}
% =====================================================================

Proposition~\ref{prop:e-representation} handles two objectives. The general case is
due to Fiaschi and Cococcioni~\cite{fiaschi2022nonarchimedean} and is reported
in~\cite{fiaschi2024informed} as follows.

\begin{theorem}[Non-standard scalarisation of an LMOP]
\label{thm:lmop-scalar}
Consider an LMOP (Definition~\ref{def:lmop}) whose priority is induced by the natural
order. Then there exists an equivalent scalar program over the same domain, whose
objective function is non-standard and has the form
\begin{equation}
  \label{eq:beta-scalarisation}
  f(x) = \beta_1 f_1(x) + \dots + \beta_n f_n(x),
\end{equation}
where $\beta_1,\dots,\beta_n \in \mathbb{E}$ are such that $\beta_{i+1}/\beta_i$ is
infinitesimal for $i = 1,\dots,n-1$.\footnote{The symbols $\beta_i$ follow the
notation of the source and are local to this theorem; they are unrelated to the real
exponent $\beta$ of the gradient-scaling rules in Chapter~\ref{ch:methods}.}
\end{theorem}

The canonical instantiation is $\beta_i = \eta^{\,i-1}$, giving
\begin{equation}
  \label{eq:eta-scalarisation}
  f(x) = \sum_{i=1}^{n} f_i(x)\,\eta^{\,i-1},
\end{equation}
which recovers Proposition~\ref{prop:e-representation} exactly when $n=2$. In the
language of Section~\ref{sec:euclidean-structure}, \eqref{eq:eta-scalarisation} is
nothing other than the canonical form \eqref{eq:canonical-form} of a Euclidean number
whose $i$-th coefficient is the $i$-th objective value: the objectives occupy
successive, strictly decreasing magnitudes, and comparison of two such numbers is
therefore the lexicographic comparison of the objective vectors.

The equivalence is conditional: it holds \emph{provided the domain is standard},
that is, provided the decision variables remain real~\cite{fiaschi2024informed}. The
hypothesis looks innocuous and is not, as Section~\ref{sec:nsa-bridge} shows the
moment those decision variables are network weights.

The use of numerical infinities and infinitesimals in optimisation is a recent but
active line of research: non-Archimedean formulations exist for linear
programming~\cite{cococcioni2018lexicographic},
mixed-integer linear programming~\cite{cococcioni2020solving}, quadratic
programming~\cite{fiaschi2022nonarchimedean}, and non-linear/evolutionary
optimisation~\cite{lai2021solving}; a survey of the parallel \emph{grossone}
methodology of Sergeyev, which pursues similar goals through a different
foundational route, is given in~\cite{sergeyev2017numerical}.

% =====================================================================
\section{Recovering a standard update direction}
\label{sec:nsa-bridge}
% =====================================================================

Theorem~\ref{thm:lmop-scalar} delivers a single objective, differentiable in each of
its standard components, so it is tempting to conclude that a hierarchical network
can simply be trained on \eqref{eq:eta-scalarisation} by gradient descent. It cannot.
Seeing exactly why is what makes the operator of Chapter~\ref{ch:methods} the
resolution of a specific difficulty rather than an unmotivated piece of linear
algebra.

\subsection{The non-standard gradient}
\label{ssec:nsa-gradient}

Let $\theta \in \mathbb{R}^{P}$ be the parameter vector of a model exposing $n$
differentiable level objectives
$\ell_1,\dots,\ell_n : \mathbb{R}^{P} \to \mathbb{R}$ ordered by decreasing priority,
and let
\begin{equation}
  \label{eq:ns-objective}
  f(\theta) \;=\; \sum_{i=1}^{n} \ell_i(\theta)\,\eta^{\,i-1} \;\in\; \mathbb{E}
\end{equation}
be the scalarisation \eqref{eq:eta-scalarisation}. Two observations make the gradient
of \eqref{eq:ns-objective} unproblematic to write down.

First, $f$ is a \emph{finite} $\mathbb{E}$-linear combination of standard real
functions. Differentiating it therefore requires no non-standard differential
calculus: the coefficients $\eta^{\,i-1}$ are constants of the field, so by
$\mathbb{E}$-linearity of the derivative,
\begin{equation}
  \label{eq:ns-gradient}
  \nabla_\theta f(\theta)
  \;=\;
  \sum_{i=1}^{n} \eta^{\,i-1}\,\nabla_\theta \ell_i(\theta)
  \;=\;
  g_1 + \eta\,g_2 + \dots + \eta^{\,n-1} g_n
  \;\in\; \mathbb{E}^{P},
  \qquad
  g_i := \nabla_\theta \ell_i(\theta) \in \mathbb{R}^{P}.
\end{equation}
Each $g_i$ is an ordinary real gradient, produced by ordinary backpropagation, and
Theorem~\ref{thm:transfer} is what guarantees that the rules used to compute it stay
valid when read in $\mathbb{E}$. Beyond that, the only operation performed is
extension of scalars on $\mathbb{R}^{P}$.

Second, \eqref{eq:ns-gradient} is already in canonical form, coordinate by
coordinate: its $j$-th component is the Euclidean number
$\sum_{i} (g_i)_j\,\alpha^{-(i-1)}$, a sum of at most $n$ monosemia with strictly
decreasing exponents. By Theorem~\ref{thm:canonical-form}, comparing two such
components is comparing the tuples $\big((g_1)_j,\dots,(g_n)_j\big)$
lexicographically. The priority structure is therefore carried by the gradient
itself, not only by the objective, which is what makes \eqref{eq:ns-gradient} the
natural object to try to descend.

\subsection{Why that gradient cannot be used directly}
\label{ssec:nsa-why-not}

\paragraph{The step leaves the standard domain.}
Theorem~\ref{thm:lmop-scalar} holds only while the decision variables remain real, as
Section~\ref{sec:lexico} flagged. A gradient step on \eqref{eq:ns-objective},
\begin{equation}
  \label{eq:ns-step}
  \theta^{+}
  \;=\; \theta - \mu\,\nabla_\theta f(\theta)
  \;=\; \theta - \mu g_1 - \mu\eta\,g_2 - \dots - \mu\eta^{\,n-1} g_n ,
\end{equation}
violates that hypothesis after a single iteration, for any real step size $\mu>0$,
whenever $g_2,\dots,g_n$ are not all zero: $\theta^{+} \in \mathbb{E}^{P}$ and
$\theta^{+} \notin \mathbb{R}^{P}$. The next forward pass would then have to be
evaluated at a non-standard argument, the level losses would themselves become
$\mathbb{E}$-valued, and the $\eta$-degrees would compound from iteration to
iteration. Nothing in this is inconsistent --- Remark~\ref{rem:algorithmic} records
that such numbers are machine-representable, and \eqref{eq:ns-step} is exactly the
iteration a non-Archimedean solver would run --- but the object being trained would
no longer be the real-weighted network whose parameters, gradients and predictions
this thesis reports.

\paragraph{The standard part discards the hierarchy.}
The canonical way back from $\mathbb{E}$ to $\mathbb{R}$ is the standard part
(Theorem~\ref{thm:standard-part}). Applied componentwise to \eqref{eq:ns-gradient},
\begin{equation}
  \label{eq:st-collapse}
  \mathrm{st}\big(\nabla_\theta f(\theta)\big)
  \;=\;
  \mathrm{st}\!\left(g_1 + \sum_{i \ge 2}\eta^{\,i-1} g_i\right)
  \;=\;
  g_1 ,
\end{equation}
because every term with $i \ge 2$ is infinitesimal and $\mathrm{st}$ retains only the
monosemium of exponent $0$. The resulting update is the gradient of the
highest-priority objective alone. This is not a lossy approximation of the
lexicographic problem: it is the sequence of programs of Definition~\ref{def:lmop}
truncated after the first, and it discards every objective the hierarchy was
introduced to express. The standard part is behaving exactly as it should; the
trouble is that the lower levels are, by construction, below the resolution of
anything it can return.

\paragraph{What is needed instead.}
Equations \eqref{eq:ns-gradient} and \eqref{eq:st-collapse} bracket the situation.
The non-standard gradient carries the full priority structure but cannot be applied
to a real parameter vector; its standard part can be applied but carries only the
first level. What is needed is a map
\begin{equation}
  \label{eq:assembler}
  \mathcal{A} : (\mathbb{R}^{P})^{n} \longrightarrow \mathbb{R}^{P},
  \qquad (g_1,\dots,g_n) \longmapsto d ,
\end{equation}
returning a \emph{real} direction $d$ in which the lower-priority objectives still
act, while the effect they may have on the higher-priority ones is controlled. Any
such map is necessarily a weakening. By Proposition~\ref{prop:debreu} no real-valued
aggregate encodes the lexicographic order, and $d$ is a real vector; so the question
is not whether the encoding is exact --- it cannot be --- but which property of the
lexicographic problem survives it.

\subsection{The property retained: first-order non-interference}
\label{ssec:nsa-what-retained}

The device used in this thesis keeps every lower-priority gradient but removes from
it the component acting along the directions already claimed by the higher-priority
ones. For two objectives,
\begin{equation}
  \label{eq:nsa-projection-sketch}
  d \;=\; g_1 + \Pi(g_2;g_1),
  \qquad
  \Pi(g_2;g_1) \;:=\; g_2 - \frac{\langle g_2,g_1\rangle}{\langle g_1,g_1\rangle}\,g_1 ,
\end{equation}
which is real, so $\theta - \mu d$ remains in $\mathbb{R}^{P}$ and the standard-domain
hypothesis of Theorem~\ref{thm:lmop-scalar} is preserved.\footnote{The quotient is
undefined when $g_1=0$; Chapter~\ref{ch:methods} skips the projection in that case
and in the numerically ill-conditioned neighbourhood of it.} What this buys follows
from $\langle g_1, \Pi(g_2;g_1)\rangle = 0$:
\begin{equation}
  \label{eq:nsa-noninterference}
  \langle g_1, d \rangle \;=\; \lVert g_1 \rVert^{2} \;\ge\; 0
  \qquad \text{for every } g_2 ,
\end{equation}
and a first-order expansion of the dominant objective along the step
$\theta^{+} = \theta - \mu d$ gives
$\ell_1(\theta^{+}) = \ell_1(\theta) - \mu\lVert g_1\rVert^{2} + o(\mu)$. To first
order the subordinate objective cannot degrade the dominant one, whatever it is and
however large it is --- which is the operational content of the lexicographic
priority, and exactly what Corollary~\ref{cor:compensation} showed no positive real
weighting can deliver. Unlike \eqref{eq:st-collapse}, the second objective survives:
it acts, undiminished, on the orthogonal complement of the first.

Three qualifications bound every claim made later in this thesis.

\begin{enumerate}
  \item The guarantee is \emph{first-order and local}. It concerns the directional
        derivative at the current $\theta$, hence the change predicted over an
        infinitesimally small step. It says nothing about the loss observed after a
        step of realistic size, nor about the trajectory.
  \item It is \emph{not} the nested optimisation of Definition~\ref{def:lmop}. The
        LMOP restricts objective $i$ to the argmin set of objectives $1,\dots,i-1$;
        \eqref{eq:nsa-projection-sketch} restricts it to the orthogonal complement of
        the higher-priority \emph{gradients} at a single point. The second is the
        first-order shadow of the first, and no more than that.
  \item With three or more objectives the composition of successive projections does
        \emph{not} preserve \eqref{eq:nsa-noninterference} in general.
        Chapter~\ref{ch:methods} exhibits a counterexample and identifies the
        parameter blocks on which the guarantee does hold.
\end{enumerate}

The construction leaves a gap, then, between the exact formulation of
Section~\ref{sec:lexico} and the real operator sketched here. That gap is what the
experiments of Chapter~\ref{ch:experiments} probe: whether enforcing the first-order
surrogate produces, in practice, the hierarchical behaviour the exact formulation
prescribes.

Note finally that $\eta$ has disappeared. It appears in \eqref{eq:ns-objective} and
\eqref{eq:ns-gradient}, where it establishes what a real substitute has to preserve,
and it occurs nowhere in the resulting update \eqref{eq:nsa-projection-sketch} --- nor
in any line of the training loop. This is the sense of
Remark~\ref{rem:algorithmic}: every non-standard quantity introduced by this chapter
is eliminated analytically before any arithmetic reaches hardware.

% =====================================================================
\section{Summary and notation}
\label{sec:nsa-summary}
% =====================================================================

\begin{itemize}
  \item Infinitesimals were the original language of the calculus, were eliminated by
        the $\varepsilon$--$\delta$ rigour of the nineteenth century, and were
        restored to rigorous standing by Robinson in the
        1960s~\cite{robinson1996nonstandard}; Alpha Theory provides an elementary
        axiomatic route to the same objects
        (Section~\ref{sec:nsa-origins}).
  \item No real-valued aggregate can encode a lexicographic priority
        (Proposition~\ref{prop:debreu}); an $\eta$-weighted sum does, faithfully
        (Proposition~\ref{prop:e-representation}), and the general $n$-objective
        statement is Theorem~\ref{thm:lmop-scalar}.
  \item A totally ordered field violating the Axiom of Archimedes contains infinite
        and infinitesimal elements; $\mathbb{R}$ cannot be such a field
        (Remark~\ref{rem:R-archimedean}).
  \item Alpha Theory postulates such a field $\mathbb{E} \supseteq \mathbb{R}$, an
        infinite number $\alpha$, and an $\alpha$-limit operator
        (Axioms~\ref{ax:field}--\ref{ax:alphalimit}). Every element of $\mathbb{E}$
        is a series of monosemia with strictly decreasing exponents
        (Theorem~\ref{thm:canonical-form}), which is what makes the ordering of
        $\mathbb{E}$ lexicographic in the coefficients.
  \item The finite elements form a set $\mathbb{E}_{\mathrm{fin}}$ strictly between
        $\mathbb{R}$ and $\mathbb{E}$, and every finite $\xi$ splits uniquely as
        $r+\zeta$ with $r$ real and $\zeta$ infinitesimal
        (Definition~\ref{def:magnitudes}, \eqref{eq:finite-decomposition}).
  \item Real functions extend to $\mathbb{E}$ (Theorem~\ref{thm:star-extension}), and
        properties whose quantifiers are all bounded by named sets transfer across
        the extension (Definition~\ref{def:elementary},
        Theorem~\ref{thm:transfer}). Transfer supplies the witnesses from the extended
        sets, which is why $\mathbb{E}$ can satisfy the transferred Archimedean
        statement and still be non-Archimedean (Remark~\ref{rem:transfer-scope}).
  \item The standard part returns the unique real number infinitely close to a finite
        element (Theorem~\ref{thm:standard-part}).
  \item Differentiating the non-standard objective gives
        $\nabla_\theta f = \sum_i \eta^{\,i-1} g_i$ with each $g_i$ an ordinary real
        gradient \eqref{eq:ns-gradient}. It cannot be descended directly: a step
        leaves $\mathbb{R}^{P}$ and voids the standard-domain hypothesis of
        Theorem~\ref{thm:lmop-scalar} \eqref{eq:ns-step}, while its standard part
        collapses to $g_1$ and deletes every level below the first
        \eqref{eq:st-collapse}.
  \item The device this thesis uses instead is a projection assembling a real
        direction $d$ in which lower-priority gradients act only orthogonally to
        higher-priority ones \eqref{eq:nsa-projection-sketch}. It retains
        first-order non-interference, $\langle g_1,d\rangle = \lVert g_1\rVert^{2}$
        \eqref{eq:nsa-noninterference}, and nothing stronger; it is not the nested
        optimisation of Definition~\ref{def:lmop}
        (Section~\ref{ssec:nsa-what-retained}).
  \item Alpha-Theoretic numbers are, in principle, machine-representable
        (Remark~\ref{rem:algorithmic}), but the construction of this thesis never
        needs to represent one: $\eta$ is eliminated analytically before any
        arithmetic reaches hardware.
\end{itemize}

\begin{table}[htbp]
  \centering
  \caption{Notation used in this chapter and in the rest of the thesis.}
  \label{tab:nsa-notation}
  \begin{tabular}{ll}
    \toprule
    \textbf{Symbol} & \textbf{Meaning} \\
    \midrule
    $\mathbb{E}$ & field of $\alpha$-Euclidean (non-standard) numbers, $\mathbb{E} \supseteq \mathbb{R}$ \\
    $\mathbb{E}_{\mathrm{fin}}$ & finite (i.e.\ non-infinite) elements of $\mathbb{E}$ \\
    $\alpha$ & the infinite number $\mathrm{num}(\mathbb{N})$ \\
    $\eta$ & the infinitesimal $\alpha^{-1}$ \\
    $\zeta$ & a generic infinitesimal, as in $\xi = r + \zeta$ \\
    $r\alpha^{k}$ & monosemium with coefficient $r$ and exponent $k$ \\
    $\lim_{n \uparrow \alpha}$ & $\alpha$-limit \\
    $\lim_{n \to \alpha}$ & ordinary (Cauchy) limit \\
    $f^{*}$ & star-extension of $f : \mathbb{R} \to \mathbb{R}$ to $\mathbb{E}$ \\
    $\mathrm{st}(\cdot)$ & standard part \\
    $x \sim y$ & $x - y$ is infinitesimal \\
    $\le_{\mathrm{lex}}$ & lexicographic order \\
    \midrule
    $\theta \in \mathbb{R}^{P}$ & parameter vector of the model \\
    $\ell_i$ & level objective of priority $i$ (lower index, higher priority) \\
    $g_i := \nabla_\theta \ell_i$ & real gradient of the $i$-th level objective \\
    $d$ & assembled real update direction \\
    $\mathcal{A}$ & assembler $(g_1,\dots,g_n) \mapsto d$ of \eqref{eq:assembler} \\
    $\Pi(g;r)$ & $g$ with its component along $r$ removed \\
    $\mu$ & step size of an idealised gradient step \\
    \bottomrule
  \end{tabular}
\end{table}

\begin{remark}[Symbols reserved by this chapter]
\label{rem:notation-scope}
No real quantity is ever denoted by $\alpha$ or $\eta$. In particular, the step size
of a gradient step is written $\mu$, per-level gradient scaling factors $\kappa_l$,
and regularisation coefficients $\lambda$ (objective-level) or $\varepsilon$
(numerical stabilisation).
\end{remark}
